% --- Archon physics formalization source begin ---
% archon:physics
% archon:covers PhyXMiniProblems/problem_phyx_mini_0673.lean
% archon:source-report reports/phyx_mini/problem_phyx_mini_0673.source.json

\chapter{Physics problem phyx\_mini\_0673}
\label{ch:PhyXMiniProblems_problem_phyx_mini_0673}

\paragraph{Problem source.}
\#\# Physical scenario

Two objects, A and B, are connected by hinges to a rigid rod that has a length \textbackslash{}( L \textbackslash{}). The objects slide along perpendicular guide rails as shown in figure. Assume object A slides to the left with a constant speed \textbackslash{}( v \textbackslash{}).

\#\# Question

Find the velocity \textbackslash{}( v\_\textbackslash{}text\{B\} \textbackslash{}) of object B as a function of the angle \textbackslash{}( \textbackslash{}theta \textbackslash{}).

\#\# Answer choices

- A: A: \textbackslash{}left(\textbackslash{}frac\{1\}\{\textbackslash{}sin \textbackslash{}theta\}\textbackslash{}right) v

- B: B: \textbackslash{}left(\textbackslash{}frac\{1\}\{\textbackslash{}cot \textbackslash{}theta\}\textbackslash{}right) v

- C: C: \textbackslash{}left(\textbackslash{}frac\{1\}\{\textbackslash{}tan \textbackslash{}theta\}\textbackslash{}right) v

- D: D: \textbackslash{}left(\textbackslash{}frac\{1\}\{\textbackslash{}sin\textbackslash{}theta\}\textbackslash{}right) v

\#\# Figure caption (auxiliary; use the image as primary evidence)

The image depicts a diagram with two points labeled \textbackslash{}( A \textbackslash{}) and \textbackslash{}( B \textbackslash{}). 

- Point \textbackslash{}( A \textbackslash{}) is positioned at the lower right, and a red vector arrow labeled \textbackslash{}( \textbackslash{}vec\{v\} \textbackslash{}) points leftward from \textbackslash{}( A \textbackslash{}). 
- Point \textbackslash{}( B \textbackslash{}) is located at the upper left.
- Points \textbackslash{}( A \textbackslash{}) and \textbackslash{}( B \textbackslash{}) are connected by a line labeled \textbackslash{}( L \textbackslash{}).
- A right triangle is formed by the line \textbackslash{}( L \textbackslash{}), the vertical distance \textbackslash{}( y \textbackslash{}), and the horizontal distance \textbackslash{}( x \textbackslash{}).
- The angle \textbackslash{}( \textbackslash{}theta \textbackslash{}) is marked between \textbackslash{}( L \textbackslash{}) and the horizontal line at point \textbackslash{}( A \textbackslash{}).
- The origin point \textbackslash{}( O \textbackslash{}) is at the intersection of two gray lines forming a right angle, representing the axes.
- The axes are labeled \textbackslash{}( x \textbackslash{}) (horizontal) and \textbackslash{}( y \textbackslash{}) (vertical).
- Distances \textbackslash{}( x \textbackslash{}) and \textbackslash{}( y \textbackslash{}) are shown as horizontal and vertical distances from \textbackslash{}( O \textbackslash{}) to \textbackslash{}( A \textbackslash{}) and \textbackslash{}( B \textbackslash{}), respectively. 

The scene represents a classical physics setup involving vectors, angles, and distances.

\#\# Dataset metadata

Kinematics; Spatial Relation Reasoning, Multi-Formula Reasoning

\paragraph{Recorded answer/context.}
C: C: \textbackslash{}left(\textbackslash{}frac\{1\}\{\textbackslash{}tan \textbackslash{}theta\}\textbackslash{}right) v

\paragraph{Figure/image path.}
/root/proposal\_for\_physic/science-mango/phyx\_mini\_run/phyx\_data/test\_image/673.png

\paragraph{Formalization target.}
create a compiling Lean file with sorry bodies at `PhyXMiniProblems/problem\_phyx\_mini\_0673.lean`. The Lean declarations must preserve the physical quantities, dimensions or dimensional roles, figure labels, governing-law hypotheses, and final relation expressed by this problem.
Use Mathlib/Physlib names found through LeanExplore where available. If a physics API is missing, introduce faithful local abstractions rather than scalar placeholder aliases.

\begin{theorem}[Physics formalization target]
\label{thm:physics:phyx_mini_0673:target}
\lean{PhyXMiniProblems.ProblemPhyXMini0673.objectB_velocity_as_function_of_angle}
\uses{def:physics:phyx-mini-0673:phyxminiproblems-problemphyxmini0673-speedinmeterspersecond, def:physics:phyx-mini-0673:phyxminiproblems-problemphyxmini0673-signedvelocityinmeterspersecond, def:physics:phyx-mini-0673:phyxminiproblems-problemphyxmini0673-perpendicularsliderrodsetup, def:physics:phyx-mini-0673:phyxminiproblems-problemphyxmini0673-figureobservation, def:physics:phyx-mini-0673:phyxminiproblems-problemphyxmini0673-hasphysicalsliderrodparameters, def:physics:phyx-mini-0673:phyxminiproblems-problemphyxmini0673-matchessuppliedfigure, def:physics:phyx-mini-0673:phyxminiproblems-problemphyxmini0673-obeysrigidrodkinematics}
The assigned autoformalize agent should translate this physics problem into Lean declarations in the covered file, with theorem and lemma proof bodies written as `by sorry`.
\end{theorem}
\begin{proof}
This is an autoformalization task, not a proof task. Produce statements that can later be proved by the physics prover without weakening the physical model.
\end{proof}

% --- Archon declaration topology begin ---
\section{Declaration topology}
\begin{definition}[Length Quantity]
  \label{def:physics:phyx-mini-0673:phyxminiproblems-problemphyxmini0673-lengthquantity}
  \lean{PhyXMiniProblems.ProblemPhyXMini0673.LengthQuantity}
  A nonnegative, unit-independent physical length
\end{definition}
\begin{proof}
  This is the declared model definition of Length Quantity.
\end{proof}
\begin{definition}[velocity Dimension]
  \label{def:physics:phyx-mini-0673:phyxminiproblems-problemphyxmini0673-velocitydimension}
  \lean{PhyXMiniProblems.ProblemPhyXMini0673.velocityDimension}
  The physical dimension of a signed velocity component
\end{definition}
\begin{proof}
  This is the declared model definition of velocity Dimension.
\end{proof}
\begin{definition}[Signed Velocity Quantity]
  \label{def:physics:phyx-mini-0673:phyxminiproblems-problemphyxmini0673-signedvelocityquantity}
  \lean{PhyXMiniProblems.ProblemPhyXMini0673.SignedVelocityQuantity}
  \uses{def:physics:phyx-mini-0673:phyxminiproblems-problemphyxmini0673-velocitydimension}
  A signed, unit-independent velocity component along one guide rail
\end{definition}
\begin{proof}
  This is the declared model definition of Signed Velocity Quantity.
\end{proof}
\begin{definition}[length In Meters]
  \label{def:physics:phyx-mini-0673:phyxminiproblems-problemphyxmini0673-lengthinmeters}
  \lean{PhyXMiniProblems.ProblemPhyXMini0673.lengthInMeters}
  \uses{def:physics:phyx-mini-0673:phyxminiproblems-problemphyxmini0673-lengthquantity}
  Read a physical length in metres
\end{definition}
\begin{proof}
  This is the declared model definition of length In Meters.
\end{proof}
\begin{definition}[speed In Meters Per Second]
  \label{def:physics:phyx-mini-0673:phyxminiproblems-problemphyxmini0673-speedinmeterspersecond}
  \lean{PhyXMiniProblems.ProblemPhyXMini0673.speedInMetersPerSecond}
  Read a nonnegative physical speed in metres per second
\end{definition}
\begin{proof}
  This is the declared model definition of speed In Meters Per Second.
\end{proof}
\begin{definition}[signed Velocity In Meters Per Second]
  \label{def:physics:phyx-mini-0673:phyxminiproblems-problemphyxmini0673-signedvelocityinmeterspersecond}
  \lean{PhyXMiniProblems.ProblemPhyXMini0673.signedVelocityInMetersPerSecond}
  \uses{def:physics:phyx-mini-0673:phyxminiproblems-problemphyxmini0673-signedvelocityquantity}
  Read a signed velocity component in metres per second
\end{definition}
\begin{proof}
  This is the declared model definition of signed Velocity In Meters Per Second.
\end{proof}
\begin{definition}[Slider Label]
  \label{def:physics:phyx-mini-0673:phyxminiproblems-problemphyxmini0673-sliderlabel}
  \lean{PhyXMiniProblems.ProblemPhyXMini0673.SliderLabel}
  The two hinged objects labelled in the supplied figure
\end{definition}
\begin{proof}
  This is the declared model definition of Slider Label.
\end{proof}
\begin{definition}[Guide Rail]
  \label{def:physics:phyx-mini-0673:phyxminiproblems-problemphyxmini0673-guiderail}
  \lean{PhyXMiniProblems.ProblemPhyXMini0673.GuideRail}
  The perpendicular guide rails shown through the point labelled O
\end{definition}
\begin{proof}
  This is the declared model definition of Guide Rail.
\end{proof}
\begin{definition}[guide Rail Of]
  \label{def:physics:phyx-mini-0673:phyxminiproblems-problemphyxmini0673-guiderailof}
  \lean{PhyXMiniProblems.ProblemPhyXMini0673.guideRailOf}
  \uses{def:physics:phyx-mini-0673:phyxminiproblems-problemphyxmini0673-sliderlabel, def:physics:phyx-mini-0673:phyxminiproblems-problemphyxmini0673-guiderail}
  Object A belongs to the horizontal rail and object B to the vertical rail
\end{definition}
\begin{proof}
  This is the declared model definition of guide Rail Of.
\end{proof}
\begin{definition}[Diagram Plane]
  \label{def:physics:phyx-mini-0673:phyxminiproblems-problemphyxmini0673-diagramplane}
  \lean{PhyXMiniProblems.ProblemPhyXMini0673.DiagramPlane}
  The Euclidean plane used by the side-view diagram, with coordinates in metres
\end{definition}
\begin{proof}
  This is the declared model definition of Diagram Plane.
\end{proof}
\begin{definition}[plane Point]
  \label{def:physics:phyx-mini-0673:phyxminiproblems-problemphyxmini0673-planepoint}
  \lean{PhyXMiniProblems.ProblemPhyXMini0673.planePoint}
  \uses{def:physics:phyx-mini-0673:phyxminiproblems-problemphyxmini0673-diagramplane}
  Construct a diagram point from its horizontal and vertical metre readouts
\end{definition}
\begin{proof}
  This is the declared model definition of plane Point.
\end{proof}
\begin{definition}[diagram Origin]
  \label{def:physics:phyx-mini-0673:phyxminiproblems-problemphyxmini0673-diagramorigin}
  \lean{PhyXMiniProblems.ProblemPhyXMini0673.diagramOrigin}
  \uses{def:physics:phyx-mini-0673:phyxminiproblems-problemphyxmini0673-diagramplane, def:physics:phyx-mini-0673:phyxminiproblems-problemphyxmini0673-planepoint}
  The rail intersection labelled O in the figure
\end{definition}
\begin{proof}
  This is the declared model definition of diagram Origin.
\end{proof}
\begin{definition}[Perpendicular Slider Rod Setup]
  \label{def:physics:phyx-mini-0673:phyxminiproblems-problemphyxmini0673-perpendicularsliderrodsetup}
  \lean{PhyXMiniProblems.ProblemPhyXMini0673.PerpendicularSliderRodSetup}
  \uses{def:physics:phyx-mini-0673:phyxminiproblems-problemphyxmini0673-lengthquantity, def:physics:phyx-mini-0673:phyxminiproblems-problemphyxmini0673-signedvelocityquantity}
  The coordinate curves are scalar readouts, not aliases for physical length: xCoordinateInMeters t and yCoordinateInMeters t are the two figure-derived components at the time readout t. The allowed motion times form a separate set so that the idealized constant-speed motion need only hold while the rod remains between the rails
\end{definition}
\begin{proof}
  This is the declared model definition of Perpendicular Slider Rod Setup.
\end{proof}
\begin{definition}[Figure Observation]
  \label{def:physics:phyx-mini-0673:phyxminiproblems-problemphyxmini0673-figureobservation}
  \lean{PhyXMiniProblems.ProblemPhyXMini0673.FigureObservation}
  The instantaneous time and angle read from the pictured configuration
\end{definition}
\begin{proof}
  This is the declared model definition of Figure Observation.
\end{proof}
\begin{definition}[slider Position]
  \label{def:physics:phyx-mini-0673:phyxminiproblems-problemphyxmini0673-sliderposition}
  \lean{PhyXMiniProblems.ProblemPhyXMini0673.sliderPosition}
  \uses{def:physics:phyx-mini-0673:phyxminiproblems-problemphyxmini0673-sliderlabel, def:physics:phyx-mini-0673:phyxminiproblems-problemphyxmini0673-diagramplane, def:physics:phyx-mini-0673:phyxminiproblems-problemphyxmini0673-planepoint, def:physics:phyx-mini-0673:phyxminiproblems-problemphyxmini0673-perpendicularsliderrodsetup}
  Position of either hinge in the perpendicular-rail coordinate frame
\end{definition}
\begin{proof}
  This is the declared model definition of slider Position.
\end{proof}
\begin{definition}[Has Physical Slider Rod Parameters]
  \label{def:physics:phyx-mini-0673:phyxminiproblems-problemphyxmini0673-hasphysicalsliderrodparameters}
  \lean{PhyXMiniProblems.ProblemPhyXMini0673.HasPhysicalSliderRodParameters}
  \uses{def:physics:phyx-mini-0673:phyxminiproblems-problemphyxmini0673-lengthinmeters, def:physics:phyx-mini-0673:phyxminiproblems-problemphyxmini0673-speedinmeterspersecond, def:physics:phyx-mini-0673:phyxminiproblems-problemphyxmini0673-perpendicularsliderrodsetup, def:physics:phyx-mini-0673:phyxminiproblems-problemphyxmini0673-figureobservation}
  Positivity and time-domain conditions for a nondegenerate physical setup
\end{definition}
\begin{proof}
  This is the declared model definition of Has Physical Slider Rod Parameters.
\end{proof}
\begin{definition}[Matches Supplied Figure]
  \label{def:physics:phyx-mini-0673:phyxminiproblems-problemphyxmini0673-matchessuppliedfigure}
  \lean{PhyXMiniProblems.ProblemPhyXMini0673.MatchesSuppliedFigure}
  \uses{def:physics:phyx-mini-0673:phyxminiproblems-problemphyxmini0673-lengthinmeters, def:physics:phyx-mini-0673:phyxminiproblems-problemphyxmini0673-perpendicularsliderrodsetup, def:physics:phyx-mini-0673:phyxminiproblems-problemphyxmini0673-figureobservation}
  Primary-image geometry at the observation time. Since θ is drawn between the rod and the leftward horizontal ray from A, x = L cos θ and y = L sin θ are the adjacent and opposite projections of the rod
\end{definition}
\begin{proof}
  This is the declared model definition of Matches Supplied Figure.
\end{proof}
\begin{definition}[Obeys Rigid Rod Kinematics]
  \label{def:physics:phyx-mini-0673:phyxminiproblems-problemphyxmini0673-obeysrigidrodkinematics}
  \lean{PhyXMiniProblems.ProblemPhyXMini0673.ObeysRigidRodKinematics}
  \uses{def:physics:phyx-mini-0673:phyxminiproblems-problemphyxmini0673-lengthinmeters, def:physics:phyx-mini-0673:phyxminiproblems-problemphyxmini0673-speedinmeterspersecond, def:physics:phyx-mini-0673:phyxminiproblems-problemphyxmini0673-signedvelocityinmeterspersecond, def:physics:phyx-mini-0673:phyxminiproblems-problemphyxmini0673-perpendicularsliderrodsetup, def:physics:phyx-mini-0673:phyxminiproblems-problemphyxmini0673-figureobservation, def:physics:phyx-mini-0673:phyxminiproblems-problemphyxmini0673-sliderposition}
  The governing physical laws. fixedHingeSeparation states rigidity directly as constant Euclidean distance between the two hinge positions. The next field formalizes A's constant leftward motion, hence the minus sign in its horizontal derivative. The final field is only the definition of the signed velocity readout of B as the derivative of its vertical coordinate
\end{definition}
\begin{proof}
  This is the declared model definition of Obeys Rigid Rod Kinematics.
\end{proof}
\begin{theorem}[squared coordinate constraint of fixed rod]
  \label{thm:physics:phyx-mini-0673:phyxminiproblems-problemphyxmini0673-squared-coordinate-constraint-of-fixed-rod}
  \lean{PhyXMiniProblems.ProblemPhyXMini0673.squared_coordinate_constraint_of_fixed_rod}
  \uses{def:physics:phyx-mini-0673:phyxminiproblems-problemphyxmini0673-lengthinmeters, def:physics:phyx-mini-0673:phyxminiproblems-problemphyxmini0673-perpendicularsliderrodsetup, def:physics:phyx-mini-0673:phyxminiproblems-problemphyxmini0673-sliderposition}
  Fixed Euclidean hinge separation gives the familiar right-triangle constraint
\end{theorem}
\begin{proof}
  The conclusion follows from the governing relations and figure readouts named in the statement of squared coordinate constraint of fixed rod.
\end{proof}
\begin{theorem}[instantaneous rigid rod velocity constraint]
  \label{thm:physics:phyx-mini-0673:phyxminiproblems-problemphyxmini0673-instantaneous-rigid-rod-velocity-constraint}
  \lean{PhyXMiniProblems.ProblemPhyXMini0673.instantaneous_rigid_rod_velocity_constraint}
  \uses{def:physics:phyx-mini-0673:phyxminiproblems-problemphyxmini0673-speedinmeterspersecond, def:physics:phyx-mini-0673:phyxminiproblems-problemphyxmini0673-signedvelocityinmeterspersecond, def:physics:phyx-mini-0673:phyxminiproblems-problemphyxmini0673-perpendicularsliderrodsetup, def:physics:phyx-mini-0673:phyxminiproblems-problemphyxmini0673-figureobservation, def:physics:phyx-mini-0673:phyxminiproblems-problemphyxmini0673-hasphysicalsliderrodparameters, def:physics:phyx-mini-0673:phyxminiproblems-problemphyxmini0673-obeysrigidrodkinematics}
  Differentiating the fixed-length constraint at the observation time gives x * x' + y * y' = 0. Here x' = -v and y' = v\_B by the two velocity laws
\end{theorem}
\begin{proof}
  The conclusion follows from the governing relations and figure readouts named in the statement of instantaneous rigid rod velocity constraint.
\end{proof}
% --- Archon declaration topology end ---
% --- Archon physics formalization source end ---
